\chapter{Verzeichnis der benutzen Hilfsmittel}

%\todo{IW: ?!}
%offenbar mit wenig Erfolg. Ich rate Ihnen ein ordentliches Rechtschreibprogramm. Word ist deutlich besser!
%Nun haben Sie ja Latex benutzt. Da scheint das nicht ganz so einfach zu sein, aber wertvolle Tipps habe ich hiergefunden:https://texwelt.de/fragen/144/wie-kann-ich-die-rechtschreibung-meines-latex-dokuments-prufen-lassen
%Wenn man ChatGPT fehlerhafte Fragen stellt (so wie diese hier), dann gibt ChatGPT auch fehlerhafte Antworten. Genau das wurde in der diesjährigen KI-Vorlesung von Prof. Annuth thematisiert.
%Es ist absolut idiotisch, ChatGPT zur Rechtschreibprüfung zu benutzen. Den Grund kann man hier sehen: Ich bekomme selten eine so fehlerhafte Version eingereicht.
%\todo{IW: Am besten raus nehmen}
\begin{itemize}
    \item \textbf{ChatGPT 4o}: Das ChatGPT-4o KI-Modell von OpenAI wurde genutzt, um den Schreibstil von Textabschnitten zu verbessern. Hierfür wurden Prompts verwendet, wie beispielsweise: \enquote{Korrigiere die Lesbarkeit des folgenden Textes. Behalte dabei die Latex-Formatierung ein.{\textbackslash}n <Text>}
\end{itemize}

%\todo{IW: Literatur}
%ist vom Stil her in Ordnung
%Wobei ich mich jetzt nicht zur Relevanz derselben äußere. Vom äußeren Anschein sieht das alles seriös aus, aber ich habe jetzt nicht die Zeit, überall hineinzusehen.
%Was mir noch unklar ist, ob Ihre oben bemängelten Ausführungen darauf beruhen, dass Sie Schwierigkeiten haben, sich geordnet auszudrücken, oder darauf, dass Sie wirklich nicht so viel Ahnung haben. Und da Sie ja ChatGPT zur Rechtschreibprüfung benutzt haben (das zeugt eher von wenig Ahnung), kann ich nicht ausschließen, dass manche Ihrer Ausführungen darauf beruhen, dass Sie das so in irgendeiner dieser Referenzen hier gelesen haben, und dann sind die eben doch nicht seriös.